\section{Prática 3}
\telaPasso{A3 · 01 — Preparar o inversor e abrir Layout XL}{39-launch-layout-xl.png}{%
\hypertarget{pratica3}{}%
\acao{Célula}{\cod{a3_inv}: um PMOS e um NMOS, W = 120 nm, L = 45 nm.}
\acao{Alimentação}{Fonte ideal no testbench. Usar as redes herdadas \cod{vdd!} e \cod{gnd!} da montagem.}
\acao{Abrir}{Launch → Layout XL}
}{O schematic usa \cod{vdd_inherit} e \cod{gnd_inherit}. Preservar esses nomes ao preparar terminais físicos e LVS, sem trocá-los silenciosamente por redes locais.}

\duasTelas{A3 · 02 — Criar a vista layout}{40-layout-startup.png}{41-nova-vista-layout.png}{%
Startup Option:\par Layout = Create New;\par Configuration = Automatic.}{%
Library: \cod{minicurso_cmos}\par Cell: \cod{a3_inv}\par View/Type: \cod{layout}\par Open with: Layout XL
}{O layout é uma nova vista da mesma célula. Confirmar os campos e clicar em OK.}

\telaAmpla{A3 · 03 — Reconhecer a área de layout}{42-layout-vazio-schematic.png}{%
À esquerda, camadas do PDK; à direita, schematic de referência.\par
A vista ainda está vazia: criar a vista não gera um layout concluído.
}{Próxima operação: gerar os dispositivos a partir do schematic no Layout XL, posicionar PMOS acima e NMOS abaixo e iniciar as conexões. As janelas intermediárias não foram registradas.}

\telaPasso{A3 · 04 — Relacionar o layout atualizado}{43-layout-inversor.png}{%
\acao{PMOS superior}{Source e corpo em \cod{vdd!}.}
\acao{NMOS inferior}{Source e corpo em \cod{gnd!}.}
\acao{vin}{Gates dos dois dispositivos.}
\acao{vout}{Drains dos dois dispositivos.}
}{Comparar com \cod{a3_inv/schematic} e conferir camadas, contatos e nomes das redes. Esta é a captura atualizada da montagem; os resultados das verificações aparecem nas próximas etapas.}

\begin{frame}{A3 · 05 — Configurar a tecnologia no Assura}
\small
\textcolor{blue_theme}{\textbf{Se a tecnologia aparecer como -undefined-}}\par\medskip
Na CIW, consultar o caminho da biblioteca:\par\smallskip
\cod{ddGetObj("gpdk045")~>readPath}\par\medskip
Em \textbf{Assura → Technology…}, selecionar:\par\smallskip
{\scriptsize\cod{/home/tools/cadence/gpdk/gpdk045_v_6_0/}\par\cod{assura_tech.lib}}\par\medskip
\textbf{Tecnologia esperada:} \cod{gpdk045_av}\par
\textbf{Rule Set:} \cod{default}
\fonte{O caminho é específico do servidor do minicurso; em outra instalação, localizar o arquivo correspondente no PDK. Reconhecer as regras é uma configuração anterior à verificação do layout.}
\end{frame}

\duasTelas{A3 · 06 — Localizar o DRC no Assura}{47-menu-assura-drc.png}{44-configuracao-drc-assura.png}{%
No layout: Assura → Run DRC…}{%
Configuração de referência:\par gpdk045\_av · default\par Regras: assuraDRC.rul
}{As regras precisam estar reconhecidas antes da execução. O próximo formulário registra a versão final usada na verificação.}

\telaPasso{A3 · 07 — Executar DRC na versão final}{48-formulario-drc-final.png}{%
\acao{Design}{\cod{a3_inv/layout}\par Área: Full}
\acao{Regras}{\cod{gpdk045_av}\par default · \cod{assuraDRC.rul}}
\acao{Executar}{Conferir nome e diretório na sessão; preservar os resultados.}
}{A captura tem Run Name vazio e Run Directory = ponto. Escolher dados de execução válidos no ambiente; o formulário sozinho não é o resultado do DRC.}

\telaPasso{A3 · 08 — Resultado registrado: DRC sem erros}{46-drc-sem-erros.png}{%
\acao{Mensagem do Assura}{\textbf{No DRC errors found.}}
\acao{O que ela confirma}{As regras de DRC executadas não encontraram violações nesta versão do layout.}
\acao{Próximo passo}{Conferir a correspondência elétrica pelo LVS.}
}{Se o layout mudar durante uma correção, repetir o DRC na versão alterada. A captura é o resultado fornecido para o minicurso.}

\duasTelas{A3 · 09 — Localizar o LVS no Assura}{49-menu-assura-lvs.png}{45-configuracao-lvs-assura.png}{%
No layout: Assura → Run LVS…}{%
Configuração de referência:\par extract.rul · compare.rul\par LVSinclude.rsf
}{Comparar schematic e layout da mesma célula \cod{a3_inv}. Tecnologia gpdk045\_av, Rule Set default.}

\telaPasso{A3 · 10 — Executar a comparação final de LVS}{50-formulario-lvs-final.png}{%
\acao{Vistas de a3\_inv}{\cod{schematic}\par\cod{layout}}
\acao{Regras}{\cod{extract.rul}\par\cod{compare.rul}\par\cod{LVSinclude.rsf}}
\acao{Executar}{Conferir nome/diretório e ler o resumo.}
}{Biblioteca \cod{minicurso_cmos}, tecnologia gpdk045\_av e Rule Set default. O resultado mostrado a seguir identifica a execução como \cod{a3_inv}.}

\telaPasso{A3 · 11 — Resultado registrado: LVS corresponde}{51-lvs-match.png}{%
\acao{Mensagem do Assura}{\textbf{Schematic and Layout Match.}}
\acao{Resumo}{Zero divergências de redes, dispositivos, pinos e parâmetros.}
\acao{Navegação}{Yes permite visualizar os resultados da nova execução.}
}{A janela também indica Total DRC violations: 0. A pergunta de navegação não é uma falha de LVS; a correspondência já está informada no resumo.}

\begin{frame}{A3 · 12 — Concluir com as evidências de DRC e LVS}
\figura{\begin{tikzpicture}[node distance=12mm,font=\small]
\node[draw=blue_theme,align=center,minimum width=2.5cm,minimum height=1.2cm] (a){\textbf{DRC}\\Sem erros};
\node[draw=blue_theme,right=of a,align=center,minimum width=2.5cm,minimum height=1.2cm] (b){\textbf{LVS}\\Schematic/Layout Match};
\node[draw=blue_theme,right=of b,align=center,minimum width=2.5cm,minimum height=1.2cm] (c){\textbf{Registrar}\\Preservar resultados};
\draw[seta] (a)--(b);\draw[seta] (b)--(c);
\node[below=10mm of b,font=\scriptsize,align=center] (d){Alterou o layout?\\Repetir as verificações pertinentes.};
\draw[seta] (d.west)-|(a.south);
\end{tikzpicture}}
\chamada{O layout verificado corresponde ao schematic do inversor.}
\fonte{Conclusão apoiada nas capturas de DRC e LVS fornecidas. A aula termina nessas verificações; extração de parasitas e simulação após extração ficam fora do escopo. Retomar o encerramento e o roadmap.}
\end{frame}
